% Imporditud: tools/import_eksperimendid.py
% Allikas: Eesti füüsikaolümpiaadi eksperimendi ülesannete
%   kogu (Rael Kalda, 2023), ülesanne Ü87, lahendus L87.
\setAuthor{EFO žürii}
\setRound{piirkonnavoor}
\setYear{2002}
\setNumber{G E2}
\setDifficulty{3}
\setTopic{geomeetriline optika}
\setVahendid{kaks läätse, joonlaud}
\setPunktid{15}
\setAllikas{Eksperimendi kogu Ü87, lahendus lk 70}
\setLahendusseis{toores}

\prob{Kaks läätse}
On antud 2 läätse. Teha kindlaks, kumb lääts on suurema optilise tugevusega ja mitu korda. On teada, et kehtib seos: 1/f = 1/a + 1/k, kus f on läätse fookuskaugus, a on eseme kaugus läätsest ja k --- kujutise kaugus läätsest.
\textit{Märkus:} kui kaks läätse on teineteisele väga lähedal (kokkupuutes), siis nende optilised tugevused liituvad algebraliselt.

\hint

\solu
1. Kindlaks teha, millist liiki läätsedega on tegu: kontrollida, kas lääts koondab valgust või mitte --- 1 p. 2. Kindlaks teha, kumb on tugevam: panna läätsed kokku ja vaadata kas süsteem koondab valgust või mitte --- 1 p. 3. Leida kumerläätse fookuskaugus: kas kasutada kauget valgusallikat või lähedast; teisel juhul tuleb mõõta ka kaugust valgusallikani ja arvutada läätse valemist 1/a + 1/k = 1/f fookuskaugus f . Idee --- 1 p, mõõtmised --- 1 p, valem --- 1 p ning arvutus --- 1 p. 4. Valemist D = 1/f arvutada kumerläätse optiline tugevus Dk. Valem --- 1 p, õige vastus --- 1 p. 5. Leida eespoolkirjeldatud viisil süsteemi fookuskaugus. Idee --- 1 p, teostus --- 1

p. 6. Arvutada süsteemi optiline tugevus D --- 1 p. 7. Leida seos süsteemi optilise tugevuse D ja kumer- ning nõgusläätse optiliste tugevuste Dk ja Dn vahel: D = Dk $-$ Dn ning arvutada Dn --- 2 p. 8. Mõõtevea hindamine ---2 p.
\probend
